\subsection{Triviality of fibered spaces and of group extensions}

PROPOSITION 2. --- Let X be a locally compact space in which a locally compact group H acts on the right, continuously and properly, by \((x,\xi)\mapsto x\xi\) . Assume that X/H is paracompact. Let g be a continuous representation of H in \(R^{n}\) . Then there exists a continuous mapping f of X into \(R^{n}\) such that \(f(x\xi)=f(x)+g(\xi)\) for all \(x\in X\) and \(\xi\in H\) .

One reduces immediately to the case that n = 1. Since the additive group R is isomorphic to the multiplicative group \(R_{+}^{*}\) , the proposition is then an immediate consequence of Prop. 7 of §2, No.~4.

COROLLARY. --- Let X be a locally compact space in which a finite-dimensional real vector space V operates on the right, continuously and properly, by \((x,v)\mapsto xv\) . Let \(\pi\) be the canonical mapping of X onto B = X/V. Assume that B is paracompact.

\begin{enumerate}
\def\labelenumi{\alph{enumi})}
\item
  There exists a continuous mapping \(f\) of \(\mathbf{X}\) into \(\mathbf{V}\) such that \(f(xv) = f(x) + v\) for all \(x \in \mathbf{X}\) and \(v \in \mathbf{V}\) .
\item
  If \(f\) is a mapping satisfying the conditions of a), then the mapping \(x \mapsto (\pi(x), f(x))\) is a homeomorphism of \(X\) onto \(B \times V\) .
\end{enumerate}

The assertion a) results from Prop. 2 in which g is taken to be the identity mapping of V. Let f be a mapping satisfying the conditions of a). The mapping \(x \mapsto x \cdot (-f(x))\) of X into X is continuous, and is constant on each orbit, hence is of the form \(\varphi \circ \pi\) , where \(\varphi\) is a continuous mapping of B into X; for every \(b \in B\) , \(\pi(\varphi(b)) = b\) . The mappings \(x \mapsto (\pi(x), f(x))\) of X into B × V and \((b, v) \mapsto \varphi(b) \cdot v\) of B × V into X are inverse to each other, because \(\varphi(\pi(x)) \cdot f(x) = x \cdot (-f(x)) \cdot (f(x)) = x\) , \(\pi(\varphi(b) \cdot v) = \pi(\varphi(b)) = b\) , and, if \(b = \pi(y)\) , then

\[
f (\varphi (\pi (y)) \cdot v) = f (y \cdot (- f (y)) \cdot v) = f (y) - f (y) + v = v.
\]

Since these mappings are continuous, they are homeomorphisms.

Remark. --- Let E be a finite-dimensional real affine space, T a compact space, \(\mu\) a measure on T of total mass 1, and f a continuous mapping of T into E. If an origin a in E is chosen, E becomes equipped with a vector space structure, and the integral \(\int_{\mathrm{T}} f(t) d\mu(t)\) therefore has meaning; it represents the point x of E such that

\[
x - a = \int_ {\mathrm{T}} (f (t) - a) d \mu (t).
\]

This point is independent of the choice of \(a\) . For, let \(a' \in \mathbf{E}\) and \(x' \in \mathbf{E}\) be such that \(x' - a' = \int_{\mathbf{T}} (f(t) - a') d\mu(t)\) . Then

\[
\begin{array}{c} x ^ {\prime} - a = (x ^ {\prime} - a ^ {\prime}) + (a ^ {\prime} - a) = \int_ {\mathrm{T}} \bigl (f (t) - a ^ {\prime} \bigr) d \mu (t) + \int_ {\mathrm{T}} (a ^ {\prime} - a) d \mu (t) \\ = \int_ {\mathrm{T}} \bigl (f (t) - a \bigr) d \mu (t) = x - a, \end{array}
\]

whence \(x' = x\) . We may therefore employ the symbol \(\int_{\mathrm{T}} f(t) d\mu(t)\) without specifying the choice of origin in E. If u is an affine mapping of E into another finite-dimensional affine space \(E'\) , then

\[
u \Big (\int_ {\mathrm{T}} f (t) d \mu (t) \Big) = \int_ {\mathrm{T}} u \big (f (t) \big) d \mu (t).
\]

For, E and \(E'\) may be identified with vector spaces in such a way that u becomes a linear mapping, in which case the formula is known (Ch. III, §3, No.~2, Prop. 2 and No.~3, Prop. 7).

Lemma 2. --- Let G be a compact group, \(\mu\) the normalized Haar measure of G, E a finite-dimensional real affine space, A the affine group of E, and \(\rho\) a homomorphism of G into A. Assume that, for every \(x \in \mathbf{E}\) , the mapping \(s \mapsto \rho(s)x\) of G into E is continuous. Then, for every \(x \in \mathbf{E}\) , the point

\[
x _ {0} = \int_ {\mathrm{G}} \rho (s) x d \mu (s) \in \mathrm{E}
\]

is invariant under G.

For, for every \(t \in G\) ,

\[
\rho (t) x _ {0} = \int_ {\mathrm{G}} \rho (t) \rho (s) x d \mu (s) = \int_ {\mathrm{G}} \rho (t s) x d \mu (s) = \int_ {\mathrm{G}} \rho (s) x d \mu (s) = x _ {0}.
\]

PROPOSITION 3. --- Let G be a locally compact group. Let H be a closed normal subgroup of G, isomorphic to \(\mathbf{R}^n\) and such that \(\mathbf{G} / \mathbf{H}\) is compact.

\begin{enumerate}
\def\labelenumi{\alph{enumi})}
\item
  There exists a closed subgroup \(\mathbf{L}\) of \(\mathbf{G}\) such that \(\mathbf{G}\) is the topological semi-direct product of \(\mathbf{L}\) and \(\mathbf{H}\) .
\item
  If \(\mathbf{M}\) is a compact subgroup of \(\mathbf{G}\) , there exists an element \(x \in \mathbf{H}\) such that \(x^{-1}\mathbf{M}x \subset \mathbf{L}\) .
\item
  Every compact subgroup of \(\mathbf{G}\) is contained in a maximal compact subgroup.
\item
  The maximal compact subgroups of \(\mathbf{G}\) are the subgroups that are the transforms of \(\mathbf{L}\) by the inner automorphisms of \(\mathbf{G}\) .
\end{enumerate}

Let \(\pi\) be the canonical homomorphism of \(\mathbf{G}\) onto \(\mathbf{K} = \mathbf{G} / \mathbf{H}\) . By passage to the quotient, the mapping \((s,h)\mapsto shs^{-1}\) of \(\mathbf{G}\times \mathbf{H}\) into \(\mathbf{H}\) defines a continuous mapping \((\sigma ,h)\mapsto \sigma \cdot h\) of \(\mathbf{K}\times \mathbf{H}\) into \(\mathbf{H}\) such that \(shs^{-1} = \pi (s)\cdot h\) . We shall identify \(\mathbf{H}\) with \(\mathbf{R}^n\) (and will therefore employ, as the case may be, either the multiplicative or the additive notation for the group law in \(\mathbf{H}\) ). By the Cor. of Prop. 2, there exists a continuous mapping \(f\) of \(\mathbf{G}\) into \(\mathbf{H}\) such that \(f(xh) = f(x) + h\) for \(x\in \mathbf{G}, h\in \mathbf{H}\) . For every \(x\in \mathbf{G}\) , let \(p(x) = x\cdot (-f(x))\) , which depends only on the coset of \(x\) with respect to \(\mathbf{H}\) . Set

\[
\begin{array}{r l} \mathrm{F} (x, y) & = p (x y) ^ {- 1} p (x) p (y) = f (x y) y ^ {- 1} x ^ {- 1} x (- f (x)) y (- f (y)) \\ & = f (x y) [ y ^ {- 1} (- f (x)) y ] (- f (y)) \\ & = f (x y) - \pi (y) ^ {- 1} f (x) - f (y). \end{array}\tag{1}
\]

We see that if \(F(x,y) = 0\) for all x,y in G, then \(p(G) = L\) is a subgroup of G that intersects each coset of H in one and only one point. Since p is continuous, G is then the topological semi-direct product of L and H (GT, III, §2, No.~10).

Now, for any \(h, h' \in H\) ,

\[
\begin{array}{r l} F (x h, y h ^ {\prime}) & = f (x h y h ^ {\prime}) - \pi (y) ^ {- 1} f (x h) - f (y h ^ {\prime}) \\ & = f (x h y) + h ^ {\prime} - \pi (y) ^ {- 1} f (x) - \pi (y) ^ {- 1} h - f (y) - h ^ {\prime} \\ & = f \big (x y (\pi (y) ^ {- 1} h) \big) - \pi (y) ^ {- 1} f (x) - f (y) - \pi (y) ^ {- 1} h \\ & = f (x y) - \pi (y) ^ {- 1} f (x) - f (y) = F (x, y). \end{array}
\]

Therefore F defines, by passage to quotients, a continuous mapping \(\varphi\) of \(K \times K\) into H.

On the other hand, for all \(x, y, z\) in \(\mathbf{G}\) , we have

\[
\begin{array}{r l} & {\mathrm{F} (z, x y) + \mathrm{F} (x, y) = f (z x y) - \pi (x y) ^ {- 1} f (z) - f (x y) + f (x y)} \\ & {\qquad - \pi (y) ^ {- 1} f (x) - f (y)} \\ & {\qquad = \pi (y) ^ {- 1} f (z x) - \pi (x y) ^ {- 1} f (z) - \pi (y) ^ {- 1} f (x) + f (z x y)} \\ & {\qquad - \pi (y) ^ {- 1} f (z x) - f (y)} \\ & {\qquad = \pi (y) ^ {- 1} \mathrm{F} (z, x) + \mathrm{F} (z x, y),} \end{array}
\]

therefore, for all \(x', y', z'\) in \(\mathbf{K}\) ,

\[
- \varphi (x ^ {\prime}, y ^ {\prime}) = \varphi (z ^ {\prime}, x ^ {\prime} y ^ {\prime}) - y ^ {- 1} \varphi (z ^ {\prime}, x ^ {\prime}) - \varphi (z ^ {\prime} x ^ {\prime}, y ^ {\prime}).
\]

Let us integrate with respect to \(z'\) by means of the normalized Haar measure \(\alpha\) of K. Setting \(\psi(x') = \int \varphi(z', x') d\alpha(z')\) , \(\psi\) is a continuous function on K, and (on observing that the operations of K in \(R^{n}\) respect the vector space structure of \(R^{n}\) by GT, VII, §2, No.~1, Prop. 1), one obtains

\[
- \varphi (x ^ {\prime}, y ^ {\prime}) = \psi (x ^ {\prime} y ^ {\prime}) - y ^ {- 1} \psi (x ^ {\prime}) - \psi (y ^ {\prime}).
\]

In other words, setting \(k - \psi \circ \pi\) , which is a continuous function on G,

\[
- \mathrm{F} (x, y) = k (x y) - \pi (y) ^ {- 1} k (x) - k (y).\tag{2}
\]

Comparing (1) and (2), one sees that if \(f\) is replaced by the continuous function \(f + k\) (which leaves verified the property \(f(xh) = f(x) + h\) ), \(F\) is replaced by 0 and, as we saw earlier, this completes the proof of a).

For every \(g \in \mathbf{G}\) , let \(l_g\) (resp. \(h_g\) ) be the unique element of \(\mathbf{L}\) (resp. \(\mathbf{H}\) ) such that \(g = h_g l_g\) . If \(h_1 \in \mathbf{H}\) and \(g \in \mathbf{G}\) , then

\[
g h _ {1} = h _ {g} l _ {g} h _ {1} = h _ {g} (l _ {g} h _ {1} l _ {g} ^ {- 1}) l _ {g},
\]

thus \(h_{gh_1} = h_g + l_g h_1 l_g^{-1}\) . For every \(g \in \mathbf{G}\) , let \(\psi_g\) be the mapping of \(\mathbf{H}\) into itself defined by

\[
\psi_ {g} (h _ {1}) = h _ {g} + l _ {g} h _ {1} l _ {g} ^ {- 1}.
\]

One sees that the mapping \((g,h_{1})\mapsto\psi_{g}(h_{1})\) of G×H into H is continuous and makes H a homogeneous space for G, in which the stabilizer of the origin is L. We observe moreover that when H is identified with \(R^{n}\) , \(\psi_{g}\) is an affine mapping of H into itself. This said, let M be a compact subgroup of G; by Lemma 2, there exists an \(x\in H\) such that \(\psi_{m}(x)=x\) for all \(m\in M\) . For \(y\in H\) , \(\psi_{y}\) is the translation with vector y; it follows that for every \(m\in M\) , \(\psi_{x^{-1}}\circ\psi_{m}\circ\psi_{x}\) transforms the origin of H into itself, therefore \(x^{-1}mx\in L\) . This proves that \(x^{-1}Mx\subset L\) , whence b).

Let \(L'\) be a closed subgroup of \(G\) containing \(L\) . Then \(L'\) is the topological semi-direct product of \(L\) and \(L' \cap H\) . If \(L'\) is compact, then \(L' \cap H\) is compact hence reduces to a point (GT, IV, §2, No.~2, Cor. 1 of Th. 2), therefore \(L' = L\) . This proves that \(L\) is a maximal compact subgroup of \(G\) ; the same is therefore true of the subgroups that are the transforms of \(L\) by the inner automorphisms of \(G\) . The assertions c) and d) of Prop. 3 are then immediate consequences of b).

PROPOSITION 4. --- Let G be a locally compact group and H a closed normal subgroup of G such that K = G/H is compact. Then every continuous representation u of H in R, such that \(u(s\xi s^{-1}) = u(\xi)\) for all \(\xi \in H\) and \(s \in G\) , may be extended to a continuous representation of G in R.

Let \(L = G \times R\) and let M be the set of \((\xi, -u(\xi))\) , where \(\xi\) runs over H. It is clear that M is a closed normal subgroup of L. Let \(L' = L/M\) and let \(\pi\) be the canonical mapping of L onto \(L'\) . The subgroup of L generated by M and R is \(H \times R\) , hence is closed; therefore \(\pi(\mathbf{R})\) is a closed subgroup N of \(L'\) . The restriction \(\rho\) of \(\pi\) to R is a bijective continuous representation of R onto N. Lemma 2 of Appendix 1 proves that \(\rho\) is bicontinuous. Moreover, \(L'/N\) is isomorphic to \(\mathrm{L}/(\mathrm{H} \times \mathbf{R}) = \mathrm{G}/\mathrm{H}\) , hence is compact. By Prop. 3, and taking into account the fact that N is in the center of \(L'\) , \(L'\) is the product of N with another subgroup. Therefore there exists a continuous representation of \(L'\) onto N that reduces on N to the identity mapping. Therefore there exists a continuous representation v of L onto R that is trivial on M and reduces on R to the identity mapping. For \(\xi \in H\) , one has \(v((\xi,0)) = v((\xi,-u(\xi))(e,u(\xi))) = u(\xi)\) , which completes the proof.

Lemma 3. --- Let G be a topological group generated by a compact neighborhood of e. Let H be a closed subgroup of G such that the homogeneous space G/H is compact. Then H is generated by a compact neighborhood of e in H.

Let C be a compact set such that G = CH. Enlarging C if necessary, we can suppose that C generates G and that G = \(\overset{\circ}{CH}\) . Then \(C^{2}\) is compact and is covered by the \(\overset{\circ}{Cs}\) ( \(s \in H\) ), which are open. Therefore there exist \(s_{1}, \ldots, s_{n}\) in H such that \(C^{2} \subset \overset{\circ}{Cs}_{1} \cup \cdots \cup \overset{\circ}{Cs}_{n}\) . Let \(\Gamma\) be the subgroup of H generated by the \(s_{i}\) . Then \(C^{2} \subset \mathrm{C}\Gamma\) . By induction, it follows that \(C^{n} \subset \mathrm{C}\Gamma\) for every n, therefore G = C \(\Gamma\) . Every element of H may be put in the form ab with \(a \in C\) , \(b \in \Gamma\) , whence \(a \in H\) , whence \(a \in C \cap H\) . Therefore H is generated by \(C \cap H\) and the \(s_{i}\) , that is, by a compact set.

Lemma 4. --- Let G be a connected topological group, D a totally disconnected normal subgroup of G. Then D is contained in the center of G.

For, let \(d \in \mathbf{D}\) . The image of \(\mathbf{G}\) under the continuous mapping \(x \mapsto x dx^{-1}\) is a connected subset of \(\mathbf{D}\) , hence reduces to \(\{d\}\) , which proves that \(xd = dx\) for all \(x \in \mathbf{G}\) .

PROPOSITION 5. --- Let G be a connected topological group admitting a discrete normal subgroup D such that K = G/D is compact, and such that the commutator subgroup of K is dense in K. Then D is finite and G is compact.

The group G is locally isomorphic to K (GT, III, §2, No.~6, Prop. 19), hence is locally compact; since it is connected, it is generated by a compact neighborhood of e. By Lemmas 3 and 4, D is a finitely generated abelian group, hence is isomorphic to a group \(Z^{r} \times D_{1}\) with \(D_{1}\) finite (A, VII, §4, No.~7, Th. 3). Suppose that r \textgreater{} 0. Then there exists a representation f of D onto Z. By Prop. 4, f may be extended to a continuous representation g of G in R. By passage to quotients, g defines a continuous representation \(g'\) of K in R/Z; since R/Z is abelian, the kernel of \(g'\) contains the commutator subgroup of K, therefore \(g'\) is trivial; in other words, \(g(\mathbf{G}) \subset \mathbf{Z}\) . Since G is connected, it follows that \(g(\mathbf{G}) = \{0\}\) , which is absurd since \(f(\mathbf{D}) = \mathbf{Z}\) . Thus r = 0 and D is finite. Consequently G is compact (GT, III, §4, No.~1, Cor. 2 of Prop. 2).
% semantic-fragments: legacy-input-seam
\relax{}
